% Include after amsmath, amssymb, amsthm, booktabs.
\section{Foundation PA to the ordinary PAeq proof presentation}
\label{sec:foundation-paeq-complete}

This section completes a separate logical interface: Foundation
Peano arithmetic, in the pinned ordered-semiring language, to the ordinary
$0,S,+,\times,=$ PAeq presentation used by the numerical compiler.
The result is a mathematical $I\Sigma_1$ proof on arbitrary internal
codes.  It does not merely remove a full-PA model assumption from a theorem
statement.  The two places that previously used that assumption are
replaced: formula substitution uses the repaired $I\Sigma_1$ construction
argument, and repeated proof opening/closing uses the explicit constructors
below.  The last subsections bind the result to the current Foundation
predicate used by \file{nativePAConsistency}, including the one changed
finite axiom and the metatheory-base comparison.  No correctness assertion
about the numerical compiler is assumed.

\subsection{The exact presentations and numerical syntax}

The source is the pinned Foundation theory \file{Peano}: its finite
\file{PeanoMinus} branch has 17 arithmetic and nine equality axioms,
and its second branch is unrestricted induction.  Its primitive language
is $0,1,+,\times,=,<$.  The target is the exact theory
\file{Presentation.ordinaryPAeq} of
\file{PAPresentationAxioms.lean}: six equality/congruence axioms,
six recursion/injectivity axioms, and successor induction with every free
parameter universally closed.  Addition and multiplication recurse on
their first argument.  This choice fixes the proof predicate; arbitrary
alternative axiom recognizers are not substituted for it.

The term map sends $1$ to $S0$ and otherwise preserves $0,+,\times$ and
variables.  Formula translation $\tau$ preserves truth, falsity,
conjunction, disjunction and quantifiers; it translates equality
componentwise and translates strict order by
\[
 \tau(s<t)=\exists z\,(\tau(s)+S z=\tau(t)),
\]
with the displayed binder implemented by bound-variable shift.
Negative atoms are the normal-form negations of their translated positive
atoms.  These are exactly the numerical functions \file{termCode} and
\file{formulaCode 0}; all constructors are the pinned Foundation
constructors, not the later radix-four formula codes.

Write $\operatorname{Alls}(b,k)$ for $k$ iterated universal quantifiers,
$F_k=(v_0,\ldots,v_{k-1})$ for the free-variable vector, and
$\operatorname{Open}(k,b)=\operatorname{subst}(F_k,b)$.
Let $\operatorname{sh}$ increment free-variable indices and let
$\operatorname{free}(p)=\operatorname{subst}_1(v_0,\operatorname{sh}p)$.
The condition $\operatorname{sh}b=b$ expresses absence of free variables.
All these operations and their validity predicates are primitive recursive.

The following numerical syntax identities hold uniformly in $I\Sigma_1$:
translation commutes with normal-form negation, free shift, bound support,
iterated universal closure, and simultaneous substitution of translated
terms.  The last identity is
\[
 \tau(\operatorname{subst}(w,p))
 =\operatorname{subst}(\operatorname{termVector}(w),\tau(p)) . \tag{S}
\]
Its induction is on the actual substitution-construction trace carrying
$w$ as an input, not an induction on an unbounded assertion about all
vectors.  Atomic steps use the term substitution identity and the one
fresh bound variable in strict order.  Boolean steps compose the child
identities; a quantifier passes the lifted vector
$(\#0,\operatorname{bShift}(w))$ to its child.  Translation commutes with
this vector lift.  This is the repaired argument in
\file{pa194-paeq-formula-subst/PA194PAeqFormulaSubst.lean}, including its
\file{formulaCode_subst_internal} theorem.  The primitive function
definitions agree with the explicit translation just given.  Thus (S)
applies to nonstandard formula codes as well as quoted standard formulas.

\subsection{Two explicit proof constructors}

Use the source's actual ten-rule one-sided proof code.  A constructor has
the form $\langle\!\langle s,\mathrm{tag},\ldots\rangle\!\rangle+1$,
where $s$ is its concluding finite-set sequent.  The relevant tags are
0 for initial excluded middle, 4 for universal introduction, 5 for
existential introduction, 6 for weakening, 7 for shift, 8 for cut,
and 9 for a theory axiom.  The definitions are in the pinned
\file{Bootstrapping/Syntax/Proof/Basic.lean}.  These constructors are
primitive recursive even where the arbitrary-model Lean interface marks
them \file{noncomputable}.

Let $\operatorname{Proof}_U(d,a)$ mean that $d$ proves the singleton
sequent $\{a\}$ in ordinary PAeq.  For a one-bound-variable matrix $p$,
define $\operatorname{CloseOne}(p,d)$ as follows:
\[
 \operatorname{allIntro}\bigl(\{\forall p\},p,
   \operatorname{wkRule}(\{\operatorname{free}(p)\}
        \cup\operatorname{sh}\{\forall p\},d)\bigr).       \tag{C1}
\]
The premise of the literal universal rule is exactly its indicated
weakened sequent.  Consequently $I\Sigma_1$ verifies
\[
 \operatorname{Proof}_U(d,\operatorname{free}(p))
 \ \longrightarrow\ \operatorname{Proof}_U(\operatorname{CloseOne}(p,d),\forall p).
\]

For the opposite direction, put $u=\operatorname{sh}p$,
$a=\forall u$, $g=\operatorname{subst}_1(v_0,u)=\operatorname{free}(p)$,
and $s=\{g\}$.  Define the following five proof nodes:
\begin{align*}
 d_0&=\operatorname{shiftRule}(\{a\},d),\\
 d_1&=\operatorname{wkRule}(\{a\}\cup s,d_0),\\
 d_3&=\operatorname{axL}(\{\neg g,\neg a\}\cup s,g),\\
 d_2&=\operatorname{exsIntro}(\{\neg a\}\cup s,\neg u,v_0,d_3),\\
 \operatorname{OpenOne}(p,d)&=\operatorname{cutRule}(s,a,d_1,d_2).
\end{align*}
Negation commutes with substitution, and $\neg a=\exists\neg u$,
so these are precisely legal instances of the literal rules.  Therefore
\[
 \operatorname{Proof}_U(d,\forall p)
 \ \longrightarrow\ \operatorname{Proof}_U(\operatorname{OpenOne}(p,d),\operatorname{free}(p))
                                                               \tag{O1}
\]
is proved uniformly in $I\Sigma_1$.  This derivation specifies the
numerical constructor, rather than inferring a primitive recursive
transformation just from the existential statement of provability.

\subsection{Replacing the full-PA closure induction}

Fix $k,b$ with $b$ a $k$-bound-variable semiformula and
$\operatorname{sh}b=b$.  For $j\le k$ define the single, explicitly
coded chain
\[
 C_j=\operatorname{subst}(F_j,\operatorname{Alls}(b,k-j)).
\]
Then $C_0=\operatorname{Alls}(b,k)$ and $C_k=\operatorname{Open}(k,b)$.
For $j<k$, set
\[
 q_j=\operatorname{subst}\bigl(\operatorname{qVec}(F_j),
                    \operatorname{Alls}(b,k-j-1)\bigr).
\]
The substitution equations give
\[
 C_j=\forall q_j,\qquad \operatorname{free}(q_j)=C_{j+1}. \tag{Chain}
\]
The second equality is exactly the vector identity
\file{free_opening_succ}, applied to the shift-fixed body
$\operatorname{Alls}(b,k-j-1)$.  Its precondition is rank $j+1$,
obtained by removing the displayed $k-j-1$ binders.  The vector computation
is: shift free indices in $\operatorname{qVec}(F_j)$, then substitute
$v_0$ for its first bound variable, obtaining $F_{j+1}$.
These identities are already $I\Sigma_1$ syntax facts and have no
provability premise.

Starting with $d_0=d$, iterate
$d_{j+1}=\operatorname{OpenOne}(q_j,d_j)$ for $j<k$.
This defines a PR constructor $\operatorname{OpenProof}(k,b,d)=d_k$.
Starting with a proof $e_0$ of $C_k$, iterate in reverse order
$e_{j+1}=\operatorname{CloseOne}(q_{k-j-1},e_j)$, defining
$\operatorname{CloseProof}(k,b,e_0)=e_k$.

The correctness induction for the first loop is just
$\operatorname{Proof}_U(d_j,C_j)$ for $j\le k$; for the second it is
$\operatorname{Proof}_U(e_j,C_{k-j})$.  The initial parameters $b,k,d$
are fixed.  Proof checking is $\Delta_1$, and the functions constructing
the displayed codes are provably total PR functions.  Each invariant
therefore has a $\Sigma_1$ presentation, so $I\Sigma_1$ induction applies.
Equivalently, one can carry the finite vector of constructed proof nodes
as the existential witness.  There is no unbounded universal quantifier
over bodies or proofs inside either induction predicate.

We have thus proved in $I\Sigma_1$, with explicit PR witnesses, that
\[
 \operatorname{Prov}_U(\operatorname{Alls}(b,k))
 \longleftrightarrow \operatorname{Prov}_U(\operatorname{Open}(k,b)) .\tag{OC}
\]
This replaces the full-PA induction in the pinned
\file{FoundationPAeqCodeClosureProof.provable_alls_iff_open}.
It also covers vacuous surplus closure parameters; it does not assert
that a noncanonical closure is itself a recognized induction axiom.

\subsection{The exact bounded induction witnesses}

For the source, an accepted induction axiom $c$ supplies
\[
 m,b\le c,\quad c=\operatorname{Alls}(b,m),\quad
 \operatorname{sh}b=b,\quad\operatorname{bv}(b)=m,
\]
and a unary matrix $K$ satisfying
\[
 K\le\operatorname{Open}(m,b),\qquad
 \operatorname{Open}(m,b)=\operatorname{Ind}_{+}(K).
\]
Here $\operatorname{Ind}_{+}$ uses $x+1$ in its step.
For the target, the \emph{exact} recognizer in
\file{FoundationPAeqCodeInductionRecognizer.lean} is
\begin{align*}
\operatorname{targetInductionR}(c)\ \longleftrightarrow\quad &\exists m\le c\ \exists b\le c\ [\,
  c=\operatorname{Alls}(b,m)\\
 &{}\land\operatorname{IsUFormula}(b)
   \land\operatorname{sh}b=b\land\operatorname{bv}(b)=m\\
 &{}\land\exists K\le\operatorname{Open}(m,b)\,
  (\operatorname{IsSemiformula}(1,K)\\
 &\hspace{29mm}{}\land
   \operatorname{Open}(m,b)=\operatorname{Ind}_{S}(K))\,].
\end{align*}
The bound on $K$ is the computed opened-body code, \emph{not} $c$.
Search first over $m,b\le c$, then over
$K\le\operatorname{Open}(m,b)$, checking the explicit PR relations.
This is total bounded primitive recursion, with a default failure value.
It returns witnesses on every accepted input, including internally
nonstandard accepted codes.  An annotated proof may instead carry these
same witnesses at each induction-axiom leaf.

\subsection{A uniform supplier for the original induction branch}

Take source witnesses $m,b,K$ as above, and set $B=\tau(b)$,
$k=\tau(K)$.  The syntax identities, including (S), give
\[
 \tau(c)=\operatorname{Alls}(B,m),\quad
 \operatorname{sh}B=B,\quad\operatorname{bv}(B)=m,\quad
 \operatorname{Open}(m,B)=\operatorname{Ind}_{+}(k).        \tag{W}
\]
This is the original \file{translated_induction_witness} argument
with its repaired substitution identity.  It now requires only $I\Sigma_1$.

Constructor inversion in the last identity extracts a matrix $q$ from
the final universal conclusion of $B$.  Explicitly, in normal form
$B=\alpha\vee(\beta\vee\forall q)$, and
\[
 \operatorname{IsSemiformula}(m+1,q),\quad
 \operatorname{sh}q=q,\quad
 \operatorname{subst}(\operatorname{qVec}(F_m),q)=k .       \tag{M}
\]
The extraction uses the three displayed outer constructors, not an
unbounded inverse search.  These are the already $I\Sigma_1$ facts of
\file{additive_matrix_witness}.

Put $m_0=\operatorname{bv}(\forall q)\le m$.  Define
\[
 B_0=\operatorname{IndClosed}_{S}(m_0,q)
\]
using base substitution vector $(0,\#0,\ldots,\#(m_0-1))$ and step
vector $(S\#0,\#1,\ldots,\#m_0)$.  Then $B_0$ is shift fixed and
has \emph{exactly} support $m_0$: the final $\forall q$ supplies the
lower bound, and both substituted premises have support at most $m_0$.
Consequently
$A_0=\operatorname{Alls}(B_0,m_0)$ is a recognized target induction
axiom, with its actual canonical parameter count.  Its one-node axiom
proof is $a_0=\operatorname{axm}(\{A_0\},A_0)$.

Apply $\operatorname{OpenProof}(m_0,B_0,a_0)$.  The result proves
\[
 \operatorname{Ind}_{S}\bigl(
       \operatorname{subst}(\operatorname{qVec}(F_{m_0}),q)\bigr).
\]
Substitution ignores surplus vector entries beyond the proved bound
support.  Therefore its matrix equals the matrix in (M), namely $k$.
The all-model proof is
\file{canonical_matrix_opening_eq}; its proof uses $I\Sigma_1$
structural substitution agreement, with both vector lengths explicit.

There is a uniform target proof of
$\operatorname{Ind}_{S}(k)\to\operatorname{Ind}_{+}(k)$.
Its construction uses a fixed target proof of $x+S0=Sx$, the formula
equality-replacement constructor for $k$, universal introduction, and a
fixed propositional template.  The formula replacement constructor is
obtained recursively on the actual formula code, with its atomic,
Boolean and two binder cases given in
\file{FoundationPAeqCodeEqualityFormula.lean}; its internal
\file{replace_aux} theorem uses $\Sigma_1$ structural induction.
The fixed arithmetic proof is quoted in
\file{FoundationPAeqCodeInductionStep.lean}.  These proofs already
require only $I\Sigma_1$ and do not assume source PA in the metatheory.
Composing them gives an actual target proof of $\operatorname{Ind}_{+}(k)$.
Finally use $\operatorname{CloseProof}(m,B,-)$ and (W) to prove $\tau(c)$.

Every step is a bounded witness extraction, syntax computation, fixed
proof constructor, or primitive recursion on a finite syntax/proof object.
This explicitly defines the induction supplier $\mathcal A_{\mathrm{ind}}(c)$.
For the equality-replacement constructor, recursion must not be justified
by a decrease of the raw formula \emph{code}: the binder case opens a
variable and can increase that code.  Instead form its finite call tree.
The query at a child is computed from the parent query by the displayed
variable renamings, before either child proof is constructed.  These
renamings preserve the connective/quantifier skeleton.  Descending into
a binder or Boolean child strictly decreases its constructor count;
atomic term replacement descends into the actual term tree.  All branching
degrees are at most two.  For initial syntax size $r$ the query tree has
at most $2^{r+1}-1$ nodes.  Compute these queries first, then construct
their proof outputs in reverse order.  The queries use PR syntax
operations, and every proof node uses one of finitely many explicit
constructor templates.  Bounded maxima of those operations, iterated for
$2^{r+1}$ steps, give PR bounds on the query payloads and proof outputs.
This is a PR implementation of the cited equality-replacement argument;
it does not appeal to termination of an unrestricted proof search.
The reasoning above proves its correctness in $I\Sigma_1$ for all accepted
source induction codes.  In particular the target canonical closure count
$m_0$ is not silently replaced by the source count $m$.

\subsection{The finite branch and the complete proof translation}

For each of the 26 source finite tags, the pinned
\file{finiteImageProof} is an actual closed ordinaryPAeq proof of its
translated sentence.  The arithmetic tags are add-zero, associativity,
commutativity, positive difference, zero lower bound, zero-less-one,
positive-one lower bound, addition monotonicity, the corresponding
multiplication laws, distributivity, and the three order laws.  The nine
equality tags cover reflexivity, symmetry, transitivity, both constants,
addition, multiplication, equality, and strict order.  The exact inventory
and Type-valued proof constants are in
\file{PAPresentationAxiomWitness.lean} and
\file{FoundationPAeqCodeFiniteSupplier.lean}.
The already checked $I\Sigma_1$ classifier says that the finite branch
accepts exactly one of these standard numerical codes, even in a
nonstandard model.  Select its proof by 26 code comparisons.

Combining this finite selection with $\mathcal A_{\mathrm{ind}}$ gives
a total PR axiom supplier $\mathcal A$, with a fixed default value on
nonaxioms, for which
\[
 I\Sigma_1\vdash
 \operatorname{Ax}_{\mathrm{PA}}(c)
 \longrightarrow \operatorname{Proof}_U(\mathcal A(c),\tau(c)). \tag{A}
\]

Define a proof translator $\mathcal D$ by recursion over the actual coded
source derivation.  Translate every formula in its finite-set sequent.
For tags 0 through 8, use the same literal proof constructor with
translated formula/term arguments and recursively translated premises.
The substitution, free, negation and sequent-shift identities make the
premise sequents exactly those required by the target rules.  For tag 9,
replace the axiom leaf $c$ by $\mathcal A(c)$ and weaken its singleton
sequent into the translated conclusion sequent.  Give invalid codes a
fixed default result.

To make primitive recursion explicit, compute a table of results for
codes $0,\ldots,d$.  Every immediate premise code of a valid constructor
is smaller than its parent code.  Each table cell thus uses only previous
cells and PR operations.  Induction on this bounded table, or the
existing $\Sigma_1$ derivation induction, proves
\[
 I\Sigma_1\vdash
 \operatorname{DerivationOf}_{\mathrm{PA}}(d,s)
 \longrightarrow
 \operatorname{DerivationOf}_{U}(\mathcal D(d),\tau[s]).   \tag{D}
\]
The existing ten-case proof in
\file{pa194-paeq-proof-transport/PA194PAeqProofTransport.lean}
is the corresponding conditional $I\Sigma_1$ argument; (A) now supplies
the actual PA axiom branch.  Because $\tau(\bot)=\bot$, we obtain
\begin{theorem}[Original PA to ordinary PAeq]
For the presentations fixed above there is a PR proof transformation
$\mathcal D$ such that
\[
 I\Sigma_1\vdash
 \operatorname{Proof}_{\mathrm{PA}}(d,\bot)
 \longrightarrow \operatorname{Proof}_{\mathrm{ordinaryPAeq}}(\mathcal D(d),\bot).
\]
In particular $I\Sigma_1$ proves the corresponding implication between
the two refutation-existence sentences, without a full-PA model assumption.
\end{theorem}

\subsection{Binding the current Foundation predicate}
\label{sec:current-foundation-adapter}

The root ProveIt revision pins its current Foundation submodule to
\file{32e1a0956a8622fad067328ca1959729a7634428}.
Use $\mathrm{PA}_c$ for this current theory and $\mathrm{PA}_h$ for the
historical source used above.  The following coding comparison fixes
the interface on \emph{all} natural codes, rather than only on quoted
standard sentences.

Both versions use the same square pairing
\[
 P(a,b)=
 \begin{cases}
   b^2+a,&a<b,\\
   a^2+a+b,&b\le a,
 \end{cases}
 \qquad
 \langle a_0,\ldots,a_k\rangle
   =P(a_0,\langle a_1,\ldots,a_k\rangle).
\]
Here the last two entries use $P$ itself.  Vectors have empty code $0$
and $\operatorname{cons}(a,v)=P(a,v)+1$; finite sets use binary-bit
membership.  This is the Foundation syntax pairing, distinct from
the machine's Cantor pairing.  In both versions, bound and free
variables have codes $P(0,i)+1$ and $P(1,i)+1$, and a function application
has code $\langle2,k,f,v\rangle+1$.
The symbols $(0,1,+,\times)$ have arity/code pairs
$(0,0),(0,1),(2,0),(2,1)$; equality and strict order have
pairs $(2,0),(2,1)$.

The formula tags, in order $0,\ldots,7$, are positive relation, negative
relation, truth, falsity, conjunction, disjunction, universal quantifier
and existential quantifier.  Their tuple equations are identical.
The proof constructors have the same ten tuple equations and
\file{Derivation.Phi} has the same ten clauses, including all
membership, premise-sequent, variable-scope, shift and weakening
conditions.  These direct definition comparisons justify identity on
terms, formulas and the \emph{logical part} of proof checking.
For validity predicates this identity can be proved by the bounded
table computation on codes up to the input code: every immediate
syntax or proof child has smaller code.  The table invariant is
$\Delta_1$ and the matching constructor equations give its successor
step.  Shift and substitution are identified by their actual finite
construction traces, with input vectors carried as parameters as in
(S).  Thus this comparison extends to internal codes in an
$I\Sigma_1$ model; it is not restricted to an external induction on
standard quoted formulas.
The falsity code in both versions is $P(3,0)+1=13$.
There is no claim that their theory-axiom predicates are identical.

In the current pin,
\file{Incompleteness/InductionSchemeDelta1.lean} defines
\begin{align*}
\operatorname{InductionUnivR}(c)\ \longleftrightarrow\quad &\exists m\le c\ \exists b\le c\ [\,
 c=\operatorname{Alls}(b,m)\\
 &{}\land\operatorname{IsUFormula}(b)
 \land\operatorname{sh}b=b\land\operatorname{bv}(b)=m\\
 &{}\land\exists K\le\operatorname{Open}(m,b)\,
  (\operatorname{IsSemiformula}(1,K)\\
 &\hspace{29mm}{}\land\operatorname{Open}(m,b)=\operatorname{Ind}_{+}(K))\,].
\end{align*}
Its \file{indBodyVal} is exactly
$K[0]\to((\forall(K\to K[\#0+1]))\to\forall K)$.
The historical \file{InductionR (fun _ => True)} has these same
conditions with one redundant true conjunct.  The current theorem
\file{PA_delta1Definable} uses the disjunction of this recognizer
and the finite \file{PeanoMinus} recognizer.  Its all-model defining
lemmas identify those relations in every $I\Sigma_1$ model.
Consequently the induction supplier constructed above applies to the
\emph{current} induction branch without changing its witness bounds.

\subsection{The changed finite axiom: an explicit logical adapter}

There are exactly 25 shared finite leaves: nine equality/congruence
axioms and 16 arithmetic axioms.  The remaining leaf,
\file{PeanoMinus.Axiom.addEqOfLt}, is
\begin{align*}
 \mathrm{PA}_h:\quad&
 \forall x\,\forall y\,
   (x<y\to\exists z\,(z<y+1\land x+z=y)),\\
 \mathrm{PA}_c:\quad&
 \forall x\,\forall y\,
   (x<y\to\exists z\,(x+z=y)).
\end{align*}
The old notation $\exists z<^{+}y$ expands to the first displayed
existential.  Translation $\tau$ is homomorphic for the displayed
connectives and quantifiers, so only a logical existential projection is
needed to reuse the existing ordinaryPAeq proof of the stronger formula.

Here is the exact projection constructor, valid for any unary matrices
$B,E$, including matrices containing further binders and free parameters.
Put
\[
 \phi=\neg B\vee\neg E,\quad
 s=\{\forall\phi,\exists E\},\quad
 r=\operatorname{sh}s\cup\{\operatorname{free}(\phi)\},\quad
 e=\operatorname{free}(E),\quad
 I=(\exists(B\land E)\to\exists E).
\]
Define six nodes:
\begin{align*}
 p_0&=\operatorname{axL}
       (r\cup\{e,\neg\operatorname{free}(B),\neg e\},e),\\
 p_1&=\operatorname{orIntro}
       (r\cup\{e\},\neg\operatorname{free}(B),\neg e,p_0),\\
 p_2&=\operatorname{exsIntro}(r,\operatorname{sh}E,v_0,p_1),\\
 p_3&=\operatorname{allIntro}(s,\phi,p_2),\\
 p_4&=\operatorname{wkRule}(s\cup\{I\},p_3),\\
 \operatorname{Erase}(B,E)&=
       \operatorname{orIntro}(\{I\},\forall\phi,\exists E,p_4).
\end{align*}
Every premise sequent is exactly the sequent required by its literal
constructor.  In particular
$\operatorname{free}(\phi)=
\neg\operatorname{free}(B)\vee\neg e$ and
$\operatorname{subst}_1(v_0,\operatorname{sh}E)=e$.
No theory leaf is used.

For completeness, fixed propositional composition need not hide a proof
search.  For arbitrary formulas $A,X,Y$, let
$a=A\land\neg X$, $b=X\land\neg Y$, $c=\neg A$, $e=Y$,
and $s=\{a,b,c,e\}$.
Three excluded-middle leaves on $A,X,Y$, followed by conjunction
introduction for $b$ and then $a$, prove $s$:
the three premises are respectively $s\cup\{A\}$,
$s\cup\{\neg X,X\}$, and $s\cup\{\neg X,\neg Y\}$.
One weakening and three disjunction introductions then prove
\[
 (A\to X)\to((X\to Y)\to(A\to Y)).
\]
This is a fixed nine-node constructor template.
Modus ponens is also an explicit five-new-node template: for proofs of
$P$ and $F=P\to Q$, weaken the latter into $\{F,Q\}$;
conjunction-introduce $\neg F=P\land\neg Q$ in
$\{\neg F,Q\}$, using the former proof and the excluded-middle
leaf on $Q$; then cut on $F$ to conclude $\{Q\}$.

Let $d_h$ be the existing fixed ordinaryPAeq proof of the translated
historical \file{addEqOfLt}.  Open its two outer parameters using
$\operatorname{OpenProof}(2,-,d_h)$.  Apply the projection just given,
the nine-node transitivity template and two modus-ponens instances.
Finally use $\operatorname{CloseProof}(2,-,-)$ for the weaker body.
The result is a specified finite proof $d_c$ of the translated current
leaf.  The function \file{weaken_difference_axiom} in the supplied
checker gives this literal construction.  Thus the current finite
supplier selects the old 25 shared proofs and $d_c$ by 26 comparisons.

\subsection{The metatheory-base comparison}

Using an old $I\Sigma_1$ syntax lemma also requires justifying its
arithmetic base; the changed finite leaf cannot be ignored at that level.
The current finite arithmetic base itself proves the historical bounded
leaf.  Here is a fixed derivation with no induction.

First, \file{zeroLtOne} and \file{addLtAdd} give
$0+y<1+y$.  The shared \file{addZero} and \file{addComm}
axioms rewrite this as $y<y+1$.
Now suppose $x<y$ and choose $z$ from the current
\file{addEqOfLt}, so $x+z=y$.
\file{zeroLe} says $0=x$ or $0<x$.
In the first case, equality substitution and $0+z=z$ give $z=y$,
hence $z<y+1$.  In the second case,
\file{addLtAdd} applied to $0<x$ gives $0+z<x+z$,
therefore $z<y$, and \file{ltTrans} with $y<y+1$ gives
$z<y+1$.  In both cases introduce
$\exists z\,(z<y+1\land x+z=y)$, discharge $x<y$, and
universally close $x,y$.
This uses only shared finite axioms, the current changed axiom, and
fixed first-order proof templates.  Conversely the historical bounded
leaf implies the current leaf by the six-node projection above.
The two finite arithmetic bases are therefore mutually derivable.

The historical metatheory writes $I\Sigma_1$ with a prenex
$\Sigma_1$ hierarchy, whereas the current metatheory uses
\file{Arithmetic.Hierarchy}.  In the historical
\file{FirstOrder/Syntax/Classical/BoundingHierarchy.lean},
$\operatorname{PrenexHierarchy}(\Sigma,s+1,\phi)$ unfolds to
$\phi=\exists\psi$ with
$\operatorname{PrenexHierarchy}(\Pi,s,\psi)$;
at $s=0$ the latter predicate is the bounded closure
$\mathcal B[<].\operatorname{Closure}(\psi)$.
Its eight constructors are truth, falsity, positive relation, negative
relation, conjunction, disjunction, bounded universal and bounded
existential quantification.  The bounds have the displayed
\file{Positive} side condition.
The current \file{Arithmetic/Basic/Hierarchy.lean} has those same
eight constructors at level zero with the same bound condition.
Structural recursion on the historical bounded derivation therefore
produces a current $\Pi_0$ derivation of $\psi$.  The current
\file{Hierarchy.sigma} constructor adds its outer existential to
give a current $\Sigma_1$ derivation of $\phi$.
Thus current $\Sigma_1$ induction includes every historical prenex
$\Sigma_1$ instance; the argument is an explicit structural inclusion,
not an identification of differently named hierarchy classes.
Together with the fixed base-axiom derivation, it allows all historical
$I\Sigma_1$ syntax lemmas used above to be applied in a current
$I\Sigma_1$ model.  The new inductions in this note have explicit
$\Delta_1$ or $\Sigma_1$ predicates as shown above.

\subsection{Current proof and native-refutation endpoint}

Define $\mathcal A_c$ by the current 26-leaf finite selection and
the same induction supplier $\mathcal A_{\mathrm{ind}}$.
The preceding arguments establish, in current $I\Sigma_1$,
\[
 \operatorname{Ax}_{\mathrm{PA}_c}(c)
 \ \longrightarrow\
 \operatorname{Proof}_{U}(\mathcal A_c(c),\tau(c)).
\]
Use it in the table recursion (D), on the current ten-rule proof
predicate, to obtain $\mathcal D_c$.  This construction works
directly on the current proof; it does not assert equality of the
historical and current PA axiom predicates.

\begin{theorem}[Exact current Foundation endpoint]
There is a total PR function $\mathcal D_c$ for which current
$I\Sigma_1$ proves
\[
 \operatorname{Proof}_{\mathrm{PA}_c}(d,13)
 \ \longrightarrow\
 \operatorname{Proof}_{\mathrm{ordinaryPAeq}}(\mathcal D_c(d),13).
\]
Both proof predicates are the numerical singleton-sequent predicates
with the coding and axiom recognizers fixed in this section.
The theorem holds uniformly for internally nonstandard $d$.
\end{theorem}

Finally, the exact source sentence in
\file{DiophantineConsistency/IPRAPA196Consistency.lean} is
\[
 \file{paRefutation}
 =\operatorname{provabilityPred}(\mathrm{PA}_c,\bot).
\]
The definitions map it into the PR language as
\file{mappedPARefutation}; its intuitionistic native reading is
\file{nativePARefutation}, and
\file{nativePAConsistency} is its negation.
The existing all-model
\file{PA195RefutationBridge.mapped_refutation_iff}
identifies the mapped formula, in every PR-language structure whose
arithmetic reduct satisfies current $I\Sigma_1$, with
$\exists d\,\operatorname{Proof}_{\mathrm{PA}_c}(d,13)$.
The endpoint theorem therefore gives, in exactly those structures,
\[
 \operatorname{Eval}(\file{mappedPARefutation})
 \ \longrightarrow\
 \exists e\,\operatorname{Proof}_{\mathrm{ordinaryPAeq}}(e,13).
\]
This is the source endpoint needed by the separate ordinaryPAeq-to-C4
compiler and the one-way $I\Sigma_1$-to-iPRA transfer.  This note does
not substitute a different consistency sentence for
\file{nativePAConsistency}.

The conclusion is a written mathematical proof, with explicit
constructors and their induction invariants.  The supplemental checker
replays 308 closure chains, 140 canonical induction readbacks,
36 existential projections and four complete changed-leaf adapters,
checking 8,828 distinct primitive proof nodes.  These finite checks
support the displayed constructor equations; they are not a proof
by testing over all natural codes.  No Lean kernel rebuild of the
entire current-package adapter is claimed.
